\documentclass[conference]{IEEEtran}
\IEEEoverridecommandlockouts
% The preceding line is only needed to identify funding in the first footnote. If that is unneeded, please comment it out.
\usepackage{cite}
\usepackage{amsmath,amssymb,amsfonts}
\usepackage{algorithmic}
\usepackage{graphicx}
\usepackage{textcomp}
\usepackage{xcolor}
\usepackage{url}
\usepackage[hidelinks]{hyperref}
\def\BibTeX{{\rm B\kern-.05em{\sc i\kern-.025em b}\kern-.08em
    T\kern-.1667em\lower.7ex\hbox{E}\kern-.125emX}}

\begin{document}

\title{面向一比特SAR成像的双向距离-方位向上采样
\thanks{如有资助机构请在此注明。如无，请删除。}
}

\author{\IEEEauthorblockN{1\textsuperscript{st} Given Name Surname}
\IEEEauthorblockA{\textit{dept. name of organization (of Aff.)} \\
\textit{name of organization (of Aff.)}\\
City, Country \\
email address or ORCID}
~\\
\and
\IEEEauthorblockN{2\textsuperscript{nd} Given Name Surname*}
\IEEEauthorblockA{\textit{dept. name of organization (of Aff.)} \\
\textit{name of organization (of Aff.)}\\
City, Country \\
email address or ORCID}
*通讯作者
~\\
\and
\IEEEauthorblockN{3\textsuperscript{rd} Given Name Surname}
\IEEEauthorblockA{\textit{dept. name of organization (of Aff.)} \\
\textit{name of organization (of Aff.)}\\
City, Country \\
email address or ORCID}
}

\maketitle

\begin{abstract}
一比特合成孔径雷达成像通过仅保留符号信息来降低采样精度，但由此产生的量化失真会严重退化图像质量。现有研究主要通过阈值设计或采集后恢复来改善一比特SAR成像。本文研究了一个互补的设计变量：在一比特量化之前如何分配固定的数字频域上采样预算。我们提出了双向距离-方位向上采样（BRAU），该方案在距离和方位两个维度上应用补零上采样，并保持固定的总分配预算 $Q_{\mathrm{up}}=R\times A$。因此，BRAU 不同于重建后的普通图像域插值。在七个SAR数据集上的实验表明，在所有测试的整数预算下，最佳双向分配在PSNR和SSIM上均一致优于最佳单向分配。该优势在非整数因子以及零阈值、非随机常数阈值和随机阈值量化下同样成立。频谱分析进一步表明，R2A2在距离压缩后具有最低的支撑区外能量、距离泄漏比和方位泄漏比。这些结果表明，在一比特SAR成像中，将量化前的频域上采样预算分配到距离和方位两个维度是一种有效的分配策略。
\end{abstract}

\begin{IEEEkeywords}
一比特SAR成像，上采样，距离-方位分配，量化
\end{IEEEkeywords}

\section{引言}

\begin{figure}[!t]
\centerline{\includegraphics[width=\linewidth]{../MD/2D SplitRT As0.6.jpg}}
\caption{SplitRT 设置 ($A_s=0.6$) 下使用的二维SAR场景示例。}
\label{fig:intro_scene_examples}
\end{figure}

合成孔径雷达通过对沿合成孔径采集的回波进行相干处理来形成高分辨率图像，已成为遥感和目标观测的核心感知手段 \cite{cumming2005sar,curlander1991sar}。传统SAR系统通常依赖高分辨率模数转换器来保留回波的幅度和相位信息。在星载SAR中，星上存储、功耗和下传链路带宽对原始回波数据流施加了额外压力，从而推动了原始数据压缩和块自适应量化策略的发展 \cite{kwok1989baq,attema2010fdbaq}。一比特SAR成像通过仅保留接收回波的符号来应对这一约束，这大幅降低了数据表示量，但也丢失了幅度信息并引入了严重的量化失真。

早期的符号编码SAR研究表明，在适当的处理条件下，仍然可以从仅符号数据中形成有意义的SAR图像 \cite{franceschetti1991signum}。近年来的一比特SAR方法通过设计量化阈值或重建策略进一步改善了图像质量。单频阈值被证明可以抑制一比特SAR成像中的谐波干扰 \cite{zhao2019onebit}，而时变和固定阈值策略也被研究用于提高低精度SAR回波的保真度 \cite{demir2018onebit,zhao2021lowprecision,nie2025fixed}。基于学习的反量化方法也被引入以在数据采集后抑制一比特谐波 \cite{si2023dequantization}。这些研究表明，阈值设计和后处理可以从严重量化的SAR数据中恢复有用信息。

大多数现有工作关注于如何构造阈值、如何从一比特测量中恢复图像，或如何在特定处理方向上使用一维上采样。相比之下，本文研究了一个受星载星上压缩场景启发的较少被探索的设计问题 \cite{magli2003lossy}：在一比特回波编码和下传之前，有限的数字上采样预算应如何在两个SAR维度上分配？设 $R$ 和 $A$ 分别表示距离和方位上采样因子，$Q_{\mathrm{up}}=R\times A$ 为总数字上采样分配量。单向分配使用 $R=Q_{\mathrm{up}},A=1$ 或 $R=1,A=Q_{\mathrm{up}}$，而双向分配在相同的 $Q_{\mathrm{up}}$ 下使用 $R>1$ 且 $A>1$。

本文提出双向距离-方位向上采样（BRAU）作为一种面向一比特SAR成像的固定预算分配策略。BRAU 在阈值生成和一比特量化之前对原始回波网格进行补零上采样，而非对重建后的图像进行插值。其工程意义不仅是采样网格的放大：BRAU 探索了二维过采样如何改变一比特量化失真相对于有效距离-方位SAR支撑区的频谱分布。在七个SAR数据集上的实验表明，在所有测试的整数预算下，最佳双向分配在PSNR和SSIM上均一致优于最佳单向分配。该优势在非整数上采样因子下同样成立，并在零阈值、非随机常数阈值和随机阈值设置下保持显著。机理分析进一步表明，双向分配在距离压缩后产生更低的支撑区外能量和方向性频谱泄漏，支持了距离-方位分配更好地匹配SAR回波二维频谱结构的解释。

主要贡献总结如下：
\begin{itemize}
    \item 我们提出了一比特SAR成像中固定预算的频域距离-方位上采样分配方案，其中量化前的总数字预算 $Q_{\mathrm{up}}=R\times A$ 保持不变。
    \item 我们证明，在多个测试预算下，双向分配在重建质量上一致优于最佳单向分配。
    \item 我们验证了该优势对非整数上采样因子和不同阈值构造具有鲁棒性。
    \item 我们提供了频谱泄漏证据，表明双向分配在距离压缩后产生更干净的二维频谱结构。
\end{itemize}

\section{方法}

\subsection{固定预算距离-方位向上采样}

设 $s(\tau,\eta)\in\mathbb{C}$ 表示复原始SAR回波，其中 $\tau$ 和 $\eta$ 分别为距离向时间和方位向时间维度。回波可写为 $s=I+jQ$，其中 $I,Q\in\mathbb{R}$ 表示同相和正交分量。在一比特SAR成像中，采样回波通过仅保留符号信息进行量化，可写为
\begin{equation}
    y = \operatorname{sign}(s + u),
\end{equation}
其中 $u\in\mathbb{C}$ 表示量化阈值。对于复回波，符号运算符分别应用于同相和正交通道：
\begin{equation}
    y=\operatorname{sign}(\operatorname{Re}(s+u))
    +j\operatorname{sign}(\operatorname{Im}(s+u)).
\end{equation}
这一操作大幅降低了采样精度，但也丢弃了幅度信息并引入了严重的量化失真。

本文研究在一比特量化之前应如何分配数字频域上采样预算。设 $R$ 为距离上采样因子，$A$ 为方位上采样因子。总上采样预算定义为
\begin{equation}
    Q_{\mathrm{up}} = R \times A .
\end{equation}
对于固定的 $Q_{\mathrm{up}}$，可以构造不同的分配策略。仅距离策略使用 $R=Q_{\mathrm{up}}$ 且 $A=1$，仅方位策略使用 $R=1$ 且 $A=Q_{\mathrm{up}}$。双向策略在相同总预算下使用 $R>1$ 且 $A>1$。无上采样情况 $R=1,A=1$ 作为基线。

所提出的双向距离-方位向上采样（BRAU）策略不改变SAR成像处理器本身，而是改变一比特量化前频域网格的维度分配。对于每个测试的 $(R,A)$，原始回波通过沿距离和方位进行FFT补零来实现上采样。对于复基带回波，实现方法为沿所选维度进行FFT变换，使用 \texttt{fftshift} 将频谱移至中心，在中心化频谱两侧补零，然后使用 \texttt{ifftshift} 和逆FFT得到上采样后的回波网格；二维上采样沿距离和方位依次执行相同操作。然后在上采样网格上生成阈值，并在距离压缩之前施加一比特量化。如果距离维被上采样，则距离采样频率和距离时间网格将相应更新以用于距离压缩。距离压缩结果随后通过频域下采样裁切回原始网格，所有方法使用相同的下游距离徙动校正（RCMC）、SAR成像、感兴趣区域提取和归一化。表~\ref{alg:brau} 总结了处理流水线。因此，BRAU 是一种量化前的原始回波操作，而非重建后的图像插值步骤。预算 $Q_{\mathrm{up}}$ 代表一比特回波编码前的星上数字预处理预算和采样网格扩展因子。

\begin{table}[t]
\caption{面向一比特SAR成像的BRAU处理流水线。}
\label{alg:brau}
\begin{center}
\scriptsize
\setlength{\tabcolsep}{3pt}
\begin{tabular}{c p{0.72\linewidth}}
\hline
步骤 & 操作 \\
\hline
1 & 输入原始SAR回波 $s(\tau,\eta)$，在 $Q_{\mathrm{up}}=R\times A$ 下选择距离/方位因子 $(R,A)$。 \\
2 & 沿距离和方位进行频域补零，得到上采样后的回波 $s_{\uparrow}$。 \\
3 & 在上采样网格上生成量化阈值 $u$。 \\
4 & 将上采样回波量化为 $y=\operatorname{sign}(s_{\uparrow}+u)$。 \\
5 & 使用当前网格匹配的距离参数进行距离压缩。 \\
6 & 通过频域下采样将距离压缩数据裁切回原始网格。 \\
7 & 进行RCMC、SAR成像、ROI提取和图像归一化。 \\
8 & 使用PSNR和SSIM将归一化图像与未量化的GT进行比较。 \\
\hline
\end{tabular}
\end{center}
\end{table}

\subsection{阈值设置}

我们评估三种阈值设置。零阈值直接使用 $u=0$ 进行符号量化。非随机常数阈值使用确定的常数阈值电平。随机阈值使用固定幅度 $|U|=A_s\hat{\sigma}$ 的随机相位阈值，其中 $\hat{\sigma}= \sqrt{2/\pi}\cdot\operatorname{mean}(|s(\tau,\eta)|)$ 估计信号RMS幅度，$A_s$ 是一个无量纲缩放因子，控制阈值幅度相对于信号电平的大小。对于随机阈值，我们比较两种相位构造变体：SplitRT 将阈值相位生成为独立距离相位和方位相位向量的和（可分离，$N_r+N_a$ 自由度），而 FullRT 为每个采样点分配独立的随机相位（非可分离，$N_r\times N_a$ 自由度）。这一比较用于检验阈值相位的空间结构是否影响分配结果。

除非另有说明，主固定预算实验使用 $A_s=0.6$ 的 SplitRT。该设置仅用于提供一致的主比较；阈值鲁棒性实验验证了双向优势在 ZT、NCT 和 RT 下均保持不变。

\subsection{评估协议}

主实验使用种子 2026、$A_s=0.6$ 的 SplitRT 和七个SAR数据集：Bangkok、city1、city2、SAR-figure、field、port 和 suburb 场景。这些数据集源自 Capella Open Data IEEE Data Contest 2026 数据集。\footnote{\href{https://radiantearth.github.io/stac-browser/\#/external/capella-open-data.s3.us-west-2.amazonaws.com/stac/capella-open-data-ieee-data-contest/collection.json}{Capella Open Data IEEE Data Contest 2026 数据集}} 从每个数据集中分层选取十个窗口，共得到 70 个配对测试样本。真实图像由未量化的原始回波经过相同的SAR链形成：距离压缩、RCMC、成像、ROI提取和图像归一化。每个一比特结果与归一化后的真实图像使用峰值信噪比（PSNR）和结构相似性（SSIM）进行比较。

对于每个预算 $Q_{\mathrm{up}}$，最佳双向组和最佳单向组通过其在 70 个配对样本上的平均 PSNR 来选取。报告的增益按双向减单向计算。统计显著性使用配对 Wilcoxon 符号秩检验在相同的配对样本向量上进行检验。

\section{实验与结果}

\subsection{固定预算主实验结果}

主实验评估将固定上采样预算分配到距离和方位两个维度是否比将相同预算集中在一个维度更有效。我们测试整数预算 $Q_{\mathrm{up}}=4,6,8,9,10$。对于每个 $Q_{\mathrm{up}}$，所有可用的整数因子对被分为无上采样、仅距离、仅方位和双向分配四组。

\begin{figure}[t]
\centerline{\includegraphics[width=\linewidth]{../Exp1_MainResult_Output/Exp1_MainResult_PSNR_curves.png}}
\caption{不同距离-方位分配策略下的固定预算 PSNR 主实验结果。}
\label{fig:main_psnr}
\end{figure}

\begin{figure}[t]
\centerline{\includegraphics[width=\linewidth]{../Exp1_MainResult_Output/Exp1_MainResult_SSIM_curves.png}}
\caption{不同距离-方位分配策略下的固定预算 SSIM 主实验结果。}
\label{fig:main_ssim}
\end{figure}

表~\ref{tab:main_gain} 比较了每个预算下最佳双向分配和最佳单向分配的绝对重建质量。这里"最佳"是通过 70 个配对样本上的平均 PSNR 来选取的。最后两列报告了增益，计算为双向减单向。在所有测试预算下，最佳双向组在 PSNR 和 SSIM 上均一致优于最佳单向组。PSNR 提升范围为 0.2856 dB 到 0.5072 dB，SSIM 提升范围为 0.0169 到 0.0221。配对 Wilcoxon 符号秩检验给出所有列出预算的 $p_{\mathrm{PSNR}}\leq4.34\times10^{-10}$ 和 $p_{\mathrm{SSIM}}=3.56\times10^{-13}$。

\begin{table}[t]
\caption{每个固定预算下的最佳双向和单向结果。每项报告平均 PSNR/SSIM；最佳组按平均 PSNR 选取。}
\begin{center}
\scriptsize
\setlength{\tabcolsep}{2pt}
\begin{tabular}{c c c c c}
\hline
$Q_{\mathrm{up}}$ & 最佳双向 & 最佳单向 & $\Delta$PSNR & $\Delta$SSIM \\
\hline
4  & R2A2 (26.0623/0.8026) & R4A1 (25.7767/0.7857)  & 0.2856 & 0.0169 \\
6  & R2A3 (26.5585/0.8246) & R1A6 (26.1415/0.8037)  & 0.4170 & 0.0208 \\
8  & R2A4 (26.8194/0.8363) & R8A1 (26.3618/0.8148)  & 0.4576 & 0.0215 \\
9  & R3A3 (26.9167/0.8400) & R9A1 (26.4482/0.8186)  & 0.4685 & 0.0214 \\
10 & R5A2 (26.9702/0.8427) & R10A1 (26.4630/0.8206) & 0.5072 & 0.0221 \\
\hline
\end{tabular}
\label{tab:main_gain}
\end{center}
\end{table}

图~\ref{fig:main_psnr} 和 \ref{fig:main_ssim} 显示了 PSNR 和 SSIM 曲线。无上采样基线给出了最低的重建质量。增加总预算改善了所有上采样组，但双向组形成的曲线始终高于仅距离和仅方位基线。这一结果支持了核心主张，即上采样预算的维度分配在一比特SAR成像中是重要的。

\subsection{非整数因子分解}

为验证双向优势不是整数因子分解的伪迹，我们进一步评估了由相同频域重采样和裁切过程实现的非整数分配组。测试预算包括 $Q_{\mathrm{up}}=3,4.5,5,7.5$。在所有情况下，最佳双向分配仍然优于最佳单向分配。PSNR 增益分别为 0.2723 dB、0.3437 dB、0.3952 dB 和 0.4567 dB。相应的 SSIM 增益范围为 0.0146 到 0.0221。这些结果表明 BRAU 不限于均衡或纯整数的分配模式。

\subsection{基于频谱泄漏的机理分析}

我们进一步分析为何在相同总预算下双向分配更为有效。工作假设是SAR回波具有二维频谱结构，而一比特量化在距离和方位两个维度上都引入失真。受过采样和量化噪声整形直觉的启发，更大的量化前网格可以将量化谐波和失真移离有效支撑区。单向分配仅在一个维度上增加采样密度，而双向分配将预算分布到两个维度上，减少了量化误差与有效SAR频谱支撑区之间的重叠。

对于节点矩阵 $X$，设其二维频谱能量为
\begin{equation}
S_X(k_r,k_a)=|\mathcal{F}_2\{X\}(k_r,k_a)|^2 .
\end{equation}
使用参考频谱 $X_{\mathrm{ref}}$，支撑区掩膜定义为：若 $|\mathcal{F}_2\{X_{\mathrm{ref}}\}(k_r,k_a)|\geq \tau \max |\mathcal{F}_2\{X_{\mathrm{ref}}\}|$ 则 $M(k_r,k_a)=1$，否则为 $0$，其中 $\tau=0.35$。然后我们使用
\begin{equation}
\rho_{\mathrm{off}}=\frac{\sum S_X(1-M)}{\sum S_X},
\end{equation}
其中所有求和均在对应的频谱网格上进行。方向泄漏比使用投影频谱
\begin{equation}
P_r(k_r)=\sum_{k_a}S_X(k_r,k_a), \quad
M_r(k_r)=\max_{k_a}M(k_r,k_a),
\end{equation}
\begin{equation}
\rho_r=
\frac{\sum_{k_r}P_r(k_r)[1-M_r(k_r)]}{\sum_{k_r}P_r(k_r)},
\end{equation}
和
\begin{equation}
P_a(k_a)=\sum_{k_r}S_X(k_r,k_a), \quad
M_a(k_a)=\max_{k_r}M(k_r,k_a),
\end{equation}
\begin{equation}
\rho_a=
\frac{\sum_{k_a}P_a(k_a)[1-M_a(k_a)]}{\sum_{k_a}P_a(k_a)}.
\end{equation}

在本分析中，我们使用 $Q_{\mathrm{up}}=4$ 并比较四种代表性的分配：无上采样（R1A1）、仅距离上采样（R4A1）、仅方位上采样（R1A4）和双向分配（R2A2）。图~\ref{fig:mechanism} 可视化了最终频域裁切之前的原始距离压缩频谱，因此 R4A1、R1A4 和 R2A2 所导致的不同频谱范围仍然可见。与单向组相比，R2A2 产生了更干净的二维频谱集中度。

\begin{table}[t]
\caption{距离压缩和频域裁切后的频谱泄漏指标（$Q_{\mathrm{up}}=4$）。数值越低表示频谱集中度越好。}
\begin{center}
\scriptsize
\begin{tabular}{c c c c}
\hline
分配策略 & 支撑区外能量 & 距离泄漏 & 方位泄漏 \\
\hline
R1A1 & 0.926008 & 0.471044 & 0.101785 \\
R4A1 & 0.893954 & 0.368013 & 0.070476 \\
R1A4 & 0.901018 & 0.394427 & 0.072251 \\
R2A2 & \textbf{0.892857} & \textbf{0.363750} & \textbf{0.066680} \\
\hline
\end{tabular}
\label{tab:mechanism_metrics}
\end{center}
\end{table}

表~\ref{tab:mechanism_metrics} 量化了图~\ref{fig:mechanism} 中的视觉趋势。在距离压缩和频域裁切后，R2A2 在所有上采样组中实现了最低的支撑区外能量比、距离泄漏比和方位泄漏比。由于 RCMC 是在距离-多普勒域中的相位校正步骤，它不会实质性改变这些泄漏指标所测量的频谱能量支撑区；因此，表~\ref{tab:mechanism_metrics} 报告的是 RC 阶段的指标。这些支撑区外和方向泄漏结果支持了以下解释：双向过采样有助于将一比特量化谐波和失真移离有效的二维SAR频谱支撑区。这一证据将中间频谱与最终指标联系起来：结合上述报告的 PSNR 和 SSIM 增益，这一代表性机理案例表明，双向组在中间 RC 阶段也具有最小的频谱泄漏。

\begin{figure}[t]
\centerline{\includegraphics[width=\linewidth]{../Exp2_Mechanism_Output/Exp2_Node2_RC_Spectra.png}}
\caption{不同分配策略下 $Q_{\mathrm{up}}=4$ 的距离压缩频谱。R4A1 和 R1A4 沿一个维度扩展频谱，而 R2A2 将扩展分配到距离和方位两个维度上，产生更干净的二维频谱集中度。}
\label{fig:mechanism}
\end{figure}

\subsection{阈值设计的鲁棒性}

最后，我们检验双向优势是否依赖于特定的阈值构造。首先，在代表性预算下对 $A_s \in [0,1]$ 范围内的 SplitRT 和 FullRT 进行比较。它们的 PSNR 和 SSIM 曲线几乎重叠，仅在整个测试范围内有微小差异。这表明主要效果并不取决于随机阈值生成的内部实现方式。

\begin{figure}[!t]
\centerline{\includegraphics[width=0.98\linewidth]{../Exp3B_ZT_NCT_RT_Output/Exp3B_GainBars.png}}
\caption{ZT、NCT 和 RT 阈值设置下的双向增益。}
\label{fig:threshold}
\end{figure}

其次，我们比较 ZT、NCT 和 RT。在本鲁棒性实验中，单向参考计算为每个样本在仅距离组和仅方位组上的最大值，这比选取一个均值最佳组更为严格。对于 $Q_{\mathrm{up}}=4,6,9$，双向组在所有三种阈值设置下仍然优于这个单向参考。在 ZT 下，PSNR 增益为 0.3612 dB、0.4162 dB 和 0.4318 dB。在 NCT 下，增益增加到 0.4308 dB、0.4896 dB 和 0.5341 dB。在 RT 下，增益仍然为正，分别为 0.2018 dB、0.3306 dB 和 0.3611 dB。因此，RT 改善了绝对重建质量，但它不是双向分配效果的来源。

\section{结论}

本文研究了面向一比特SAR成像的固定预算频域距离-方位上采样分配。我们不仅探讨了如何设计量化阈值或如何在一比特采集后恢复图像，还研究了固定的量化前数字预算 $Q_{\mathrm{up}}=R\times A$ 应如何在距离和方位维度之间进行分配。实验结果表明，在所有测试的整数预算下，双向分配一致优于最佳单向分配。在非整数因子下也观察到了相同的优势，表明该效果不是由特殊的整数因子分解模式引起的。

机理分析进一步支持了分配效果。在距离压缩后，双向组比仅距离和仅方位分配产生了更低的支撑区外频谱能量和更低的方向泄漏。这表明将预算分配到两个维度上减少了量化失真与有效二维SAR频谱支撑区之间的重叠。阈值实验还表明，在 ZT、NCT 和 RT 设置下该优势仍然保持，因此观察到的效果并不依赖于特定的阈值构造。

这些发现支持 BRAU 作为一种在固定数字上采样预算下改善一比特SAR成像的简单分配策略。在星载低比特下传压缩设置下，BRAU 表明将有限的星上数字上采样分配到距离和方位两个维度比集中在一个维度更为有效。未来工作将把分析扩展到更广泛的SAR场景、硬件感知实现，以及在距离-方位处理中二维量化噪声整形的更明确理论模型。

\begin{thebibliography}{00}
\bibitem{cumming2005sar}
I. G. Cumming and F. H. Wong, \textit{Digital Processing of Synthetic Aperture Radar Data: Algorithms and Implementation}. Norwood, MA, USA: Artech House, 2005.

\bibitem{curlander1991sar}
J. C. Curlander and R. N. McDonough, \textit{Synthetic Aperture Radar: Systems and Signal Processing}. New York, NY, USA: Wiley, 1991.

\bibitem{kwok1989baq}
R. Kwok and W. T. K. Johnson, ``Block adaptive quantization of Magellan SAR data,'' \textit{IEEE Transactions on Geoscience and Remote Sensing}, vol. 27, no. 4, pp. 375--383, 1989, doi: 10.1109/36.29557.

\bibitem{attema2010fdbaq}
E. Attema, C. Cafforio, M. Gottwald, P. Guccione, A. Monti Guarnieri, F. Rocca, and P. Snoeij, ``Flexible dynamic block adaptive quantization for Sentinel-1 SAR missions,'' \textit{IEEE Geoscience and Remote Sensing Letters}, vol. 7, no. 4, pp. 766--770, 2010, doi: 10.1109/LGRS.2010.2047242.

\bibitem{magli2003lossy}
E. Magli and G. Olmo, ``Lossy predictive coding of SAR raw data,'' \textit{IEEE Transactions on Geoscience and Remote Sensing}, vol. 41, no. 5, pp. 977--987, 2003, doi: 10.1109/TGRS.2003.811556.

\bibitem{franceschetti1991signum}
G. Franceschetti, V. Pascazio, and G. Schirinzi, ``Processing of signum coded SAR signal: Theory and experiments,'' \textit{IEE Proceedings F Radar and Signal Processing}, vol. 138, no. 3, pp. 192--198, 1991, doi: 10.1049/ip-f-2.1991.0025.

\bibitem{zhao2019onebit}
B. Zhao, L. Huang, and W. Bao, ``One-bit SAR imaging based on single-frequency thresholds,'' \textit{IEEE Transactions on Geoscience and Remote Sensing}, vol. 57, no. 9, pp. 7017--7032, 2019, doi: 10.1109/TGRS.2019.2910284.

\bibitem{demir2018onebit}
M. Demir and E. Ercelebi, ``One-bit compressive sensing with time-varying thresholds in synthetic aperture radar imaging,'' \textit{IET Radar, Sonar \& Navigation}, vol. 12, no. 12, pp. 1517--1526, 2018, doi: 10.1049/iet-rsn.2018.5044.

\bibitem{zhao2021lowprecision}
B. Zhao, L. Huang, and B. Jin, ``Strategy for SAR imaging quality improvement with low-precision sampled data,'' \textit{IEEE Transactions on Geoscience and Remote Sensing}, vol. 59, no. 4, pp. 3150--3160, 2021, doi: 10.1109/TGRS.2020.3014300.

\bibitem{nie2025fixed}
G. Nie, B. Zhao, Q. Liu, L. Huang, and G. Liao, ``One-bit synthetic aperture radar imaging based on fixed-threshold with slow-time fluctuations,'' \textit{IEEE Transactions on Geoscience and Remote Sensing}, vol. 63, pp. 1--15, 2025, doi: 10.1109/TGRS.2024.3519757.

\bibitem{si2023dequantization}
C. Si, B. Zhao, L. Huang, and S. Liu, ``A convolutional de-quantization network for harmonics suppression in one-bit SAR imaging,'' \textit{IEEE Transactions on Geoscience and Remote Sensing}, vol. 61, pp. 1--16, 2023, doi: 10.1109/TGRS.2023.3330530.
\end{thebibliography}

\end{document}
